\section{Uniform spatial-radius means and Cameron regulation}
\label{S:uniform-mean-regulation}

Retain the original nonzero covariance, amplitude and real spaces of
Section~\ref{S:actual-cameron-mean}. Fix
$\alpha/4<s<1/2$ and put $\beta=1/2-s>0$.
Write $A_s=C_s^2\sqrt{R_{\mathrm{sp}}}$ for the constant in
\eqref{E:cameron-fractional-product}.
Choose a real even smooth compact $\eta$, equal to one on the unit ball,
with $0\le\eta\le1$ and support in the closed radius-two ball.
For $R\ge1$ set $\eta_R(x)=\eta(x/R)$ and $m_R=\eta_R^2$.
In this section the mean acts on $E=H^{1/2}\times H^{-1/2}$.
The constant data columns belong to $F$, outside $E$.
All uniform-radius claims below concern $E$.
The actual spatial map $T_R:H_{\mathrm{sp}}\to\mathcal K(E)$
is the $E$ restriction of Proposition~\ref{P:cameron-spatial-map}.
Its literal row is
\begin{equation}\label{E:uniform-mean-row}
  \langle T_Rqz,w\rangle_E
  =\langle q,k_R(z,w)\rangle_{H_{\mathrm{sp}}},
  \qquad
  k_R(z,w)=(\eta_Rz_1)(\eta_R\Lambda^{-1}w_2).
\end{equation}
Thus the physical mask occurs once, after covariance convolution.

\subsection{Uniform multiplier and spatial tails}

\begin{lemma}[Uniform cutoff and physical mean]
  \label{L:uniform-mean-cutoff}
  For every real $\sigma$, multiplication by $\eta_R$ on $H^\sigma$
  has norm at most
  \[
    M_\sigma=2^{|\sigma|/2}(2\pi)^{-d/2}
    \int_{\R^d}|\widehat\eta(p)|\langle p\rangle^{|\sigma|}\ud p,
  \]
  independently of $R\ge1$. With $M=M_{1/2}$ and $M_0=A_sM^2$,
  \begin{equation}\label{E:uniform-spatial-mean}
    \|k_R(z,w)\|_{H_{\mathrm{sp}}}
    \le M_0\|z\|_E\|w\|_E,\qquad
    \|T_R\|_{H_{\mathrm{sp}}\to\mathcal L(E)}\le M_0 .
  \end{equation}
\end{lemma}

\begin{proof}
  Peetre's inequality bounds the weighted Fourier kernel for
  multiplication by $\eta_R$ by convolution with
  \[
    2^{|\sigma|/2}(2\pi)^{-d/2}
    |\widehat{\eta_R}(q)|\langle q\rangle^{|\sigma|}.
  \]
  Here $\widehat{\eta_R}(q)=R^d\widehat\eta(Rq)$.
  Substitution $p=Rq$ and $\langle p/R\rangle\le\langle p\rangle$
  bound its $L^1$ norm by $M_\sigma$.
  The $L^1*L^2$ inequality proves the assertion on Schwartz
  functions and density extends it to $H^\sigma$.
  The map $w_2\mapsto\Lambda^{-1}w_2$ is an isometry from
  $H^{-1/2}$ to $H^{1/2}$, and $H^{1/2}\hookrightarrow H^s$
  has norm one. Apply the cutoff bound at $\sigma=1/2$ and
  Lemma~\ref{L:cameron-fractional-products} to the two factors
  in~\eqref{E:uniform-mean-row}. Hilbert Cauchy--Schwarz then
  gives the operator bound. This is the existing completed physical
  map: the row agrees on smooth controls and hence on their completion.
  Its compact values at each fixed $R$ are retained.
  This proves Lemma~\ref{L:uniform-mean-cutoff}.
\end{proof}

\begin{lemma}[Uniform heat and frequency tails]
  \label{L:uniform-mean-spatial-tails}
  Let $C_\varepsilon$ be the actual covariance-Hilbert heat
  contraction with Fourier multiplier $e^{-\varepsilon|\xi|^2/2}$.
  There is a real orthogonal projection $P_K$ on $H_{\mathrm{sp}}$
  corresponding to the closed frequency ball $\{|\xi|\le K\}$.
  It commutes with heat. For $K\ge1$, $\varepsilon>0$ and $R\ge1$,
  \begin{align}
    \|T_R(I-C_\varepsilon)\|
    &\le M_h\varepsilon^{\beta/2},
    &M_h&=4A_sM^2\E|G_d|^\beta,\label{E:uniform-mean-heat}\\
    \|T_R(I-P_K)\|
    &\le M_fK^{-\beta},
    &M_f&=\frac{M_h}{1-e^{-1/2}},\label{E:uniform-mean-frequency}
  \end{align}
  where $G_d$ is a standard real $d$-dimensional normal vector.
  Operator norms are from $H_{\mathrm{sp}}$ to $\mathcal L(E)$.
  The projection $P_K$ is allowed to have infinite rank.
\end{lemma}

\begin{proof}
  Translations $\tau_hf(x)=f(x-h)$ extend to strongly continuous
  isometries of the actual $H_{\mathrm{sp}}$.
  Indeed, in its covariance Fourier isometry they multiply by
  $e^{-i\xi\cdot h}$; dominated convergence gives continuity
  on tests, then density and isometry give it on the completion.
  The real subspace is preserved.
  The inequalities $|e^{-it}-1|\le2\min(1,|t|^\beta)$ and
  $|\xi|^\beta\langle\xi\rangle^s\le\langle\xi\rangle^{1/2}$
  show that
  \[
    \|\tau_hf-f\|_{H^s}
    \le2\min(1,|h|^\beta)\|f\|_{H^{1/2}} .
  \]
  Use this for $f=\eta_Rz_1$ and $g=\eta_R\Lambda^{-1}w_2$.
  The product identity
  \[
    \tau_h(fg)-fg=(\tau_hf-f)\tau_hg+f(\tau_hg-g)
  \]
  holds in $H_{\mathrm{sp}}$, first for smooth factors and
  then by the continuous fractional product extension.
  Lemma~\ref{L:uniform-mean-cutoff} gives
  \begin{equation}\label{E:uniform-mean-row-translation}
    \|\tau_hk_R(z,w)-k_R(z,w)\|_{H_{\mathrm{sp}}}
    \le4A_sM^2\min(1,|h|^\beta)\|z\|_E\|w\|_E .
  \end{equation}

  The Bochner integral $\int p_\varepsilon(h)\tau_hk\ud h$
  exists, since the translation orbit is strongly continuous
  and bounded in the separable covariance Hilbert space.
  Its Fourier image is $e^{-\varepsilon|\xi|^2/2}\widehat k$:
  ordinary Gaussian integration proves this on tests and
  contraction extends it to all $k$.
  It is therefore the actual $C_\varepsilon$ of the source construction.
  Heat has variance $\varepsilon$ here; there is no additional
  Fourier normalization factor. Integrating
  \eqref{E:uniform-mean-row-translation} with
  $h=\sqrt\varepsilon G_d$ bounds
  $\|(I-C_\varepsilon)k_R(z,w)\|$ by
  $M_h\varepsilon^{\beta/2}\|z\|_E\|w\|_E$.
  Selfadjointness in~\eqref{E:uniform-mean-row} proves
  \eqref{E:uniform-mean-heat}.

  To construct $P_K$, let $V$ be the closed image of the covariance
  Fourier isometry into $L^2(\mu)$, including its fixed
  $(2\pi)^{d/2}$ factor. No surjectivity is assumed.
  Its formula extends to Schwartz functions: compact cutoff
  approximation converges in Schwartz seminorms, and temperedness
  of $\mu$ controls the Fourier error in $L^2(\mu)$.
  Choose real even radial smooth $\chi_n$, equal to one on
  the closed ball of radius $K$, zero outside radius $K+1/n$,
  and between zero and one.
  For real Schwartz $f$, the function $\chi_n\widehat f$ is
  the transform of another real Schwartz function and so lies in $V$.
  Dominated convergence against $|\widehat f|^2\mu$ gives
  $1_{\{|\xi|\le K\}}\widehat f\in V$.
  The pointwise limit includes the boundary sphere, even if
  $\mu$ charges it. Test density and contraction extend this
  assertion to all of $V$.
  Indicator multiplication is selfadjoint and idempotent on $V$;
  pulling it back defines the asserted real orthogonal $P_K$.
  All Fourier multipliers here commute.

  On $|\xi|>K$ one has
  $1-e^{-|\xi|^2/(2K^2)}\ge1-e^{-1/2}$, whence
  \[
    \|(I-P_K)k\|_{H_{\mathrm{sp}}}
    \le(1-e^{-1/2})^{-1}\|(I-C_{K^{-2}})k\|_{H_{\mathrm{sp}}}.
  \]
  Apply the row heat bound and selfadjointness of $P_K$ in
  \eqref{E:uniform-mean-row} to obtain
  \eqref{E:uniform-mean-frequency}.
  No density of $\mu$, periodic state space or finite rank was used.
  This proves Lemma~\ref{L:uniform-mean-spatial-tails}.
\end{proof}

\subsection{The time tensor and the actual regulator}

\begin{lemma}[Dimension-free Hilbert tensor extension]
  \label{L:uniform-mean-tensor}
  If $H_0,G$ are real Hilbert spaces, $2\le r<\infty$ and
  $A:H_0\to L^r(I)$ is bounded, then
  \[
    \|A\otimes I_G\|_{H_0\otimes G\to L^r(I;G)}\le\|A\|.
  \]
  The tensor extension need not be compact when $A$ is compact.
\end{lemma}

\begin{proof}
  For a finite tensor write $v=\sum_{j=1}^n f_j\otimes e_j$
  with $(e_j)$ orthonormal. Its squared norm is $\sum_j\|f_j\|^2$.
  Let $\gamma_j$ be independent standard real normals and
  $m_r=\E|\gamma_1|^r$. Finite covariance diagonalization gives
  \[
    \E\left\|\sum_j\gamma_j f_j\right\|_{H_0}^r
    \le m_r\left(\sum_j\|f_j\|_{H_0}^2\right)^{r/2}.
  \]
  In fact, if $\lambda_l$ are the covariance eigenvalues and
  $L=\sum_l\lambda_l$, the squared norm has law
  $\sum_l\lambda_l\zeta_l^2$ for independent standard $\zeta_l$.
  For $L>0$, convexity at exponent $r/2\ge1$ bounds
  $(\sum_l(\lambda_l/L)\zeta_l^2)^{r/2}$
  by $\sum_l(\lambda_l/L)|\zeta_l|^r$.
  Taking expectations proves the inequality; $L=0$ is immediate.
  Also $L=\sum_j\|f_j\|^2$.

  Choose measurable representatives of the finitely many $Af_j$.
  Scalar Gaussian moments and Tonelli give
  \begin{align*}
    m_r\int_I\left(\sum_j|Af_j(t)|^2\right)^{r/2}\ud t
    &=\E\left\|A\sum_j\gamma_jf_j\right\|_{L^r}^r\\
    &\le\|A\|^rm_r\left(\sum_j\|f_j\|^2\right)^{r/2}.
  \end{align*}
  Cancel the finite positive $m_r$ and take the $r$th root.
  Finite tensor density and Bochner completeness give the extension.
  The zero space is trivial.
  If $Af\ne0$ and $G$ has an infinite orthonormal sequence $(e_j)$,
  the images of the bounded sequence $f\otimes e_j$ have pairwise
  distance $\sqrt2\|Af\|_{L^r}$. They have no convergent subsequence,
  even when the scalar map $A$ is compact.
  Only finite auxiliary Gaussians were used in the estimate.
  This proves Lemma~\ref{L:uniform-mean-tensor}.
\end{proof}

Recall the actual supported identity $H_I=H_t\otimes H_{\mathrm{sp}}$,
the compact scalar covariance map $C_t:H_t\to L^r(I)$, $2<r<2/a$,
and its completed tensor $U=C_t\otimes I_{H_{\mathrm{sp}}}$ from
Section~\ref{S:actual-cameron-mean}.
On finite smooth tensors $U$ agrees with the earlier $\mathcal U$;
density gives equality of their completed maps.
For block applications fix the same $r$ as in
\eqref{E:cameron-block-parameters}. All $c_t,\varepsilon_J$ and
mean norms below use that one choice.
Choose a smooth-test orthonormal basis of $H_t$ and let $\pi_J$
be its first-$J$ orthogonal projection, with $\pi_0=0$.
Put
\[
  c_t=\|C_t\|,\qquad
  \varepsilon_J=\|C_t(I-\pi_J)\|_{H_t\to L^r(I)},\qquad
  \mathcal T_R=T_RU .
\]

\begin{proposition}[Uniform completed mean regulation]
  \label{P:uniform-mean-regulator}
  On $H_I$ the operators
  \[
    P_{\mathrm{ret}}=\pi_J\otimes P_K,\qquad
    Q_{J,K}=I-P_{\mathrm{ret}}
  \]
  are real orthogonal projections, and $\varepsilon_J\to0$.
  The actual mean is a strongly measurable
  $L^r(I;\mathcal K(E))$ field and satisfies, uniformly for $R\ge1$,
  \begin{equation}\label{E:uniform-completed-mean}
    \|\mathcal T_R\|\le M_0c_t,\qquad
    \|\mathcal T_RQ_{J,K}\|
    \le M_0\varepsilon_J+M_fc_tK^{-\beta}.
  \end{equation}
  These norms have domain $H_I$ and target $L^r(I;\mathcal L(E))$.
\end{proposition}

\begin{proof}
  The tensor projections are selfadjoint, idempotent and contractive
  on finite tensors and hence on their Hilbert completion.
  For the scalar tail, failure of convergence would supply unit
  $v_J$ with $\|C_t(I-\pi_J)v_J\|$ bounded below on a subsequence.
  But $(I-\pi_J)v_J$ tends weakly to zero, since its pairing with
  $h$ is bounded by $\|(I-\pi_J)h\|$.
  Compactness supplies a norm-convergent image subsequence; bounded
  linear scalar observables show its limit is zero, a contradiction.
  Lemma~\ref{L:uniform-mean-tensor} transfers this scalar tail norm
  to the Hilbert-valued time map without a dimension factor.

  For fixed $R$, $T_R$ is bounded into the separable compact-operator
  space on $E$. Composition with the actual Bochner field $U\Phi$
  gives the stated mean and its first bound.
  The exact decomposition is
  \[
    Q_{J,K}=(I-\pi_J)\otimes I+\pi_J\otimes(I-P_K).
  \]
  The first term is bounded by $M_0\varepsilon_J$.
  For the second term use the typed intertwining
  \[
    U(\pi_J\otimes(I-P_K))
    =(I_{L^r}\otimes(I-P_K))(C_t\pi_J\otimes I).
  \]
  It holds first on finite tensors, then by continuity.
  Apply~\eqref{E:uniform-mean-frequency} pointwise, followed by
  Lemma~\ref{L:uniform-mean-tensor} for $C_t\pi_J$, whose norm
  is at most $c_t$. This proves~\eqref{E:uniform-completed-mean}.
  Infinite spatial rank is allowed throughout; neither tensor
  compactness nor uniform-radius compactness has been asserted.
  This proves Proposition~\ref{P:uniform-mean-regulator}.
\end{proof}

\subsection{The given-noise block Cameron embedding}

Use $B,\delta,N,s_m,\ell_m,I_m$ and $\theta,r,p$ exactly as in
\eqref{E:cameron-block-parameters} and its preceding block partition.
Let $i_I:H_I\to H_W$ be the isometric closed inclusion.
On the full raw noise space set
\[
  P_{\mathrm{full}}=i_IP_{\mathrm{ret}}i_I^*,\qquad
  Q_{\mathrm{full}}=I-P_{\mathrm{full}} .
\]
These are genuine orthogonal projections on $H_W$.
Define the actual residual fast noise and stopped $E$ driver by
\begin{equation}\label{E:uniform-actual-residual}
  Y_m(f)=W(Q_{\mathrm{full}}T_{\delta,m}f),\qquad
  X_m(u)=g_\delta A_{m_R}(Y_m)(\min(u,\ell_m)),\qquad
  g_\delta=\delta^{(1-a)/2}.
\end{equation}
Here $T_{\delta,m}$ is the literal raw isometry
\eqref{E:raw-block-pullback}, and $A_{m_R}$ takes its $E$ restriction.
Its isometry is proved in Lemma~\ref{L:raw-block-inputs}.
The symbol $H_{\mathrm{fast}}$ below denotes this same original
covariance norm in fast coordinates, carried isometrically into $H_W$.
The $E$ little H\"older space means the zero-start closure of Lipschitz
compact-operator paths in $C^\theta$ norm. Compact operators on separable
$E$ and polygonal approximation make this closure separable.
The last interval can be empty, in which case $X_N=0$.

\begin{proposition}[Uniform-radius actual Cameron bound]
  \label{P:uniform-actual-cameron}
  The tuple $X=(X_m)_{m=0}^N$ is a centered Gaussian variable in
  the separable little H\"older compact-operator path space on $E$.
  Its actual Bochner covariance $Lh=\E[XW(h)]$, $h\in H_W$, is
  \begin{equation}\label{E:uniform-actual-covariance}
    (Lh)_m(u)=\sqrt\delta\int_0^{\min(u,\ell_m)}
    S_E(-v)\mathcal T_R(Q_{J,K}i_I^*h)(s_m+\delta v)
    S_E(v)\ud v .
  \end{equation}
  In the exact norm
  $\|x\|_{\delta,p}=(\delta\sum_m\|x_m\|_{C^\theta}^p)^{1/p}$,
  its actual Cameron embedding has the uniform-radius bound
  \begin{equation}\label{E:uniform-actual-cameron-bound}
    \|Lh\|_{\delta,p}
    \le\sqrt\delta\,C_{B,E}(2/B)^{1/p-1/r}
    [M_0\varepsilon_J+M_fc_tK^{-\beta}]\|h\|_{H_W},
    \qquad C_{B,E}\le18B^{3-1/r}.
  \end{equation}
  All spatial constants are independent of $R\ge1$.
\end{proposition}

\begin{proof}
  Orthogonality and raw pullback isometry give
  $\operatorname{Var}Y_m(f)\le\|f\|_{H_{\mathrm{fast}}}^2$.
  Gaussian transfer, Proposition~\ref{P:source-gaussian-transfer},
  therefore constructs the drivers on the given $W$, with finite
  path moments at each fixed $R,B,\delta,J,K$.
  Choose an admissible H\"older exponent strictly above $\theta$
  as in Lemma~\ref{L:cameron-actual-residual}; constant stopping
  puts the paths in the separable little $C^\theta$ closure.
  Finite-coefficient constructions and their $L^2$ limits show
  that every continuous linear coordinate is in the first chaos
  of the given $W$.

  Every literal observed row, pulled back to slow time, belongs
  to $H_I$. This is the supported-row approximation from
  Lemma~\ref{L:cameron-actual-residual}, also valid before any
  projection; approximate endpoints inside $I$ against the
  integrable time kernel and locally finite spatial covariance.
  For such a row $k$, exact orthogonality gives
  \[
    \langle Q_{\mathrm{full}}T_{\delta,m}k,h\rangle_{H_W}
    =\langle T_{\delta,m}k,i_IQ_{J,K}i_I^*h\rangle_{H_W}.
  \]
  Directions perpendicular to $H_I$ disappear from this pairing.
  Apply the completed $J=0$ identity
  \eqref{E:cameron-actual-covariance} to $Q_{J,K}i_I^*h$,
  and restrict its input and output to $E$.
  The free $F$ group preserves $E$ and restricts to $S_E$;
  the mean restricts to $\mathcal T_R$.
  This proves~\eqref{E:uniform-actual-covariance}.
  Its amplitude is exactly
  $g_\delta\delta^{a/2}=\sqrt\delta$, on the actual noise $W$.
  Arbitrary projected controls are covered by the completed mean
  identity; no smoothness or finite rank of $P_K$ is required.
  Compact operator values and strong continuity of $S_E$ and
  its adjoint make conjugation norm continuous on compact operators,
  as in Proposition~\ref{P:cameron-actual-embedding}.
  Thus the integral is Bochner measurable.
  Fixed-parameter $L^2$ covariance continuity, dense states and
  rational times identify the entire operator path.

  For explicit free constants, the Bessel-unitary Fourier matrix
  of $S_E(v)$ is the matrix displayed in that proposition.
  For $|v|\le B$, $B\ge1$, the upper off-diagonal entry is at most
  $\sqrt2 B$: use $|\sin(vk)|\le Bk$ for $k\le1$ and
  $\langle k\rangle/k\le\sqrt2$ for $k\ge1$.
  The lower entry is at most one, as are both diagonal entries;
  at zero the sine quotient has value $v$.
  Its Frobenius bound is $\sqrt{3+2B^2}\le3B$.
  Hence $G_{B,E}=\sup_{|v|\le B}\|S_E(v)\|\le3B$.
  With
  \[
    \mathcal M_h(t)=\mathcal T_R(Q_{J,K}i_I^*h)(t),\qquad
    t_m=\|\mathcal M_h\|_{L^r(I_m;\mathcal L(E))},
  \]
  the primitive increment bound and the affine time change give
  \[
    \|(Lh)_m\|_{C^\theta}
    \le\sqrt\delta\,C_{B,E}\delta^{-1/r}t_m,\qquad
    C_{B,E}=G_{B,E}^2[B^{1-1/r}+B^{1-1/r-\theta}]
    \le18B^{3-1/r}.
  \]
  Here $1-1/r>\theta$, constant stopping preserves the bound,
  and the empty block contributes zero.

  Put $n=N+1\ge2$. The finite power-mean inequality at exponent
  $r/p>1$ with nonnegative inputs $t_m^p$ gives
  \[
    \left(\sum_m t_m^p\right)^{r/p}
    \le n^{r/p-1}\sum_m t_m^r,\qquad
    \left(\sum_m t_m^p\right)^{1/p}
    \le n^{1/p-1/r}\left(\sum_m t_m^r\right)^{1/r}.
  \]
  Raising nonnegative terms uses only positive powers, including
  when $t_m=0$. The true partition gives
  $\sum_m t_m^r=\|\mathcal M_h\|_{L^r(I)}^r$ and $\delta n\le2/B$.
  The exact $\delta$ weight therefore yields
  \[
    \|Lh\|_{\delta,p}
    \le\sqrt\delta\,C_{B,E}(\delta n)^{1/p-1/r}\|\mathcal M_h\|_{L^r}
    \le\sqrt\delta\,C_{B,E}(2/B)^{1/p-1/r}
    [M_0\varepsilon_J+M_fc_tK^{-\beta}]\|h\|_{H_W}.
  \]

  Finally let $H_{\mathrm{seen}}$ be the closed span in $H_W$
  of all coefficient tests $k_\ell$ with $\ell(X)=W(k_\ell)$.
  The first-chaos Cameron identification in
  Proposition~\ref{P:cameron-actual-embedding} applies on $H_W$:
  $L$ kills $H_{\mathrm{seen}}^\perp$ and its Cameron space is
  $L(H_{\mathrm{seen}})$ with minimum raw preimage norm.
  Thus $\|Lh\|_{\mathrm{CM}}=\|P_{\mathrm{seen}}h\|_{H_W}$.
  Applying the bound to this minimum-norm representative gives
  precisely the actual Cameron embedding bound, without an
  inverse covariance or finite-dimensional density.
  This proves Proposition~\ref{P:uniform-actual-cameron}.
\end{proof}

The orthogonal $P_{\mathrm{full}}$ and $Q_{\mathrm{full}}$ Gaussian
fields are independent; residual time blocks can remain correlated.
The retained spatial head may have infinite rank.
Heat filtering alone would not give this orthogonal decomposition.
These linear bounds do not bound the expected residual path norm
uniformly on growing masks, nor perform conditioning or integration
over the infinite head. The remaining stochastic, product and
original-point upper estimates require their own proofs.
